% ===========================================================================
% Scientific Horizons — Quarterly Review of Discovery & Innovation
% Elegant multi-column scientific magazine template for AAAS / IEEE Spectrum /
% Nature Reviews style publications. Features cover page, masthead, table of
% contents, multi-column feature articles with drop caps and pull quotes,
% infographic callout boxes, department sections, and contributor bios.
%
% Compile with:  latexmk -pdf main.tex
% Edit online at: https://letx.app
% ===========================================================================

\documentclass[a4paper,11pt]{article}

% === FONT & ENCODING ========================================================
\usepackage[utf8]{inputenc}
\usepackage[T1]{fontenc}
\usepackage{tgpagella}
\usepackage{FiraSans}
\usepackage{fontawesome5}

% === GEOMETRY ===============================================================
\usepackage[margin=20mm, top=18mm, bottom=22mm, headheight=14pt, includeheadfoot]{geometry}

% === COLOR ==================================================================
\usepackage[table]{xcolor}
\definecolor{accent}{HTML}{1B3A5C}       % deep navy
\definecolor{accent2}{HTML}{2E86AB}      % teal highlight
\definecolor{graytext}{HTML}{5A5A5A}     % secondary text
\definecolor{lightrule}{HTML}{C5CDD4}    % subtle rules
\definecolor{boxbg}{HTML}{EDF2F6}        % callout box fill
\definecolor{pagebg}{HTML}{FCFCFA}       % subtle page tint
\definecolor{coveraccent}{HTML}{152C47}  % darker navy for cover
\definecolor{bioline}{HTML}{D6DDE4}      % contributor bio rules
\pagecolor{pagebg}

% === PARAGRAPH SETUP (must be 0pt for table alignment) ======================
\setlength{\parindent}{0pt}
\setlength{\parskip}{0pt}

% === GRAPHICS & DRAWING =====================================================
\usepackage{graphicx}
\usepackage{tikz}
\usetikzlibrary{calc, positioning, shadows, shapes.geometric, backgrounds}
\usepackage[most]{tcolorbox}
\tcbuselibrary{skins,breakable}
\usetikzlibrary{shadows}
\usepackage{multicol}
\usepackage{lettrine}

% === TABLES & LISTS =========================================================
\usepackage{booktabs}
\usepackage{tabularx}
\usepackage{enumitem}
\setlist{nosep, leftmargin=*, itemsep=0pt, topsep=0pt}

% === HEADINGS ===============================================================
\usepackage{titlesec}

% Department heading (e.g. RESEARCH HIGHLIGHTS)
\titleformat{\section}
  {\sffamily\bfseries\Large\color{accent}}
  {}
  {0pt}
  {}
  [\vspace{-4pt}{\color{lightrule}\hrule height 1pt}\vspace{8pt}]
\titlespacing{\section}{0pt}{18pt}{6pt}

% Article title
\titleformat{\subsection}
  {\sffamily\bfseries\LARGE\color{accent}}
  {}
  {0pt}
  {}
\titlespacing{\subsection}{0pt}{14pt}{8pt}

% Subhead within article
\titleformat{\subsubsection}
  {\sffamily\bfseries\large\color{accent2}}
  {}
  {0pt}
  {}
\titlespacing{\subsubsection}{0pt}{12pt}{4pt}

% === HEADERS & FOOTERS ======================================================
\usepackage{fancyhdr}
\fancyhf{}
\fancyhead[LE,RO]{\sffamily\footnotesize\color{graytext}Scientific Horizons}
\fancyhead[LO,RE]{\sffamily\footnotesize\color{graytext}Vol.\ 12, No.\ 3\quad·\quad Autumn 2024}
\fancyfoot[C]{\sffamily\footnotesize\color{graytext}\thepage}
\renewcommand{\headrulewidth}{0.3pt}
\renewcommand{\footrulewidth}{0pt}
\pagestyle{fancy}

% === MATH ===================================================================
\usepackage{amsmath}
\usepackage{amssymb}

% === CAPTIONS ===============================================================
\usepackage{caption}
\captionsetup{font={sf,footnotesize,color=graytext}, labelfont={bf,color=accent}, labelsep=period}

% === HYPERLINKS (load near end) =============================================
\usepackage{hyperref}
\hypersetup{
  colorlinks=true,
  linkcolor=accent2,
  urlcolor=accent2,
  citecolor=accent,
  pdfauthor={Scientific Horizons Editorial Board},
  pdftitle={Scientific Horizons --- Vol. 12, No. 3, Autumn 2024},
  pdfsubject={Scientific Magazine Quarterly},
  pdfkeywords={science, magazine, quarterly, research, materials science, genetics}
}

% ===========================================================================
% CUSTOM COMMANDS
% ===========================================================================

% Department heading on cover / TOC
\newcommand{\coverfeature}[2]{%
  {\sffamily\bfseries\color{accent2}\footnotesize\MakeUppercase{#1}}\\[2pt]
  {\sffamily\large\color{accent}#2\par}%
}

% Article byline
\newcommand{\byline}[1]{%
  \vspace{4pt}
  {\sffamily\itshape\color{graytext}#1\par}
  \vspace{2pt}
  {\color{lightrule}\hrule height 0.5pt}
  \vspace{10pt}
}

% Pull quote --- breaks out of multicol for full-width impact
\newcommand{\pullquote}[2][2]{%
  \end{multicols}
  \vspace{8pt}
  \begin{tcolorbox}[
    enhanced,
    colback=boxbg,
    colframe=accent,
    arc=0mm,
    boxrule=0pt,
    leftrule=4pt,
    top=10pt,
    bottom=10pt,
    left=14pt,
    right=14pt,
    before skip=0pt,
    after skip=0pt,
  ]
  {\large\itshape\color{accent}``#2''}
  \end{tcolorbox}
  \vspace{8pt}
  \begin{multicols}{#1}
  \noindent\ignorespaces
}

% Infographic / callout box within column flow
\newtcolorbox{infobox}[1][]{
  enhanced,
  colback=boxbg,
  colframe=accent2!30,
  arc=3mm,
  boxrule=0.5pt,
  top=8pt,
  bottom=8pt,
  left=12pt,
  right=12pt,
  before skip=10pt,
  after skip=10pt,
  fontupper=\sffamily,
  #1
}

% Contributor bio
\newcommand{\contributor}[3]{%
  \begin{minipage}[t]{\linewidth}
    \vspace{6pt}
    {\sffamily\bfseries\color{accent}#1}\par
    {\sffamily\footnotesize\color{graytext}#2}\par
    \vspace{3pt}
    {\small#3\par}
    \vspace{2pt}
    {\color{bioline}\hrule height 0.5pt}
  \end{minipage}
  \medskip
}

% Drop cap paragraph start
\newcommand{\dropcap}[2]{%
  \lettrine[lines=3, lhang=0.25, loversize=0.08, findent=3pt, nindent=0pt]{#1}{#2}%
}

% Small decorative diamond
\newcommand{\diamondrule}{%
  \vspace{4pt}
  {\color{lightrule}\hrule height 0.5pt}
  \vspace{-2pt}
  \begin{center}
  \color{accent2}\faDiamond
  \end{center}
  \vspace{-6pt}
  {\color{lightrule}\hrule height 0.5pt}
  \vspace{6pt}
}

% ===========================================================================
% DOCUMENT
% ===========================================================================

\emergencystretch=3em\relax
\begin{document}

% ====== PAGE 1 : COVER ======================================================
\thispagestyle{empty}
\newgeometry{margin=0pt}
\pagecolor{coveraccent}

\begin{tikzpicture}[remember picture, overlay]
  % Subtle dot grid pattern in header
  \foreach \x in {0,...,10} {
    \foreach \y in {0,...,6} {
      \fill[white, opacity=0.06]
        ([xshift=\x*21mm+5mm, yshift=-\y*14mm-8mm]current page.north west) circle (2.2pt);
    }
  }

  % Right-side decorative geometric element
  \foreach \r in {1,...,6} {
    \draw[white, opacity=0.08, line width=0.6pt]
      ([xshift=-18mm, yshift=-45mm]current page.north east)
      circle (\r*6mm);
  }

  % Journal title
  \node[white, align=center, inner sep=0] at ([yshift=-38mm]current page.north) {
    {\fontfamily{ppl}\fontsize{48}{54}\selectfont\bfseries SCIENTIFIC}\\[2pt]
    {\fontfamily{ppl}\fontsize{48}{54}\selectfont\bfseries HORIZONS}
  };

  % Thin decorative line
  \draw[white, opacity=0.5, line width=1pt]
    ([yshift=-62mm]current page.north west) --
    ([yshift=-62mm]current page.north east);

  % Subtitle
  \node[white, opacity=0.85, font=\sffamily\large\itshape] at ([yshift=-72mm]current page.north) {
    A Quarterly Review of Discovery \& Innovation
  };

  % Issue info
  \node[white, opacity=0.7, font=\sffamily\footnotesize] at ([yshift=-84mm]current page.north) {
    Vol.\ 12 \quad·\quad No.\ 3 \quad·\quad Autumn 2024
  };

  % ====== White content area ======
  \fill[pagebg] (current page.south west) rectangle ([yshift=120mm]current page.south east);

  % Decorative accent stripe
  \fill[accent2]
    ([yshift=120mm]current page.south west) --
    ([yshift=120mm]current page.south east) --
    ([yshift=118mm]current page.south east) --
    ([yshift=118mm]current page.south west) -- cycle;

  % Featured articles heading
  \node[accent, font=\sffamily\bfseries\footnotesize\scshape] at ([yshift=-100mm]current page.north) {
    In This Issue
  };

  % Featured article 1
  \node[align=left, text width=130mm, inner sep=0] at ([yshift=-112mm]current page.north) {
    {\sffamily\bfseries\color{accent2}\footnotesize\MakeUppercase{Feature}}\\[1pt]
    {\sffamily\fontsize{16}{20}\selectfont\bfseries\color{accent}
     The Computational Frontier: Machine Learning Meets}\\[0pt]
    {\sffamily\fontsize{16}{20}\selectfont\bfseries\color{accent}
     Quantum Materials Discovery}\\[3pt]
    {\sffamily\footnotesize\color{graytext}\itshape
     Elena Vasquez \& James Chen\ \textbar\ page 4}
  };

  % Rule
  \draw[lightrule, line width=0.4pt]
    ([yshift=-125mm]current page.north) ++(-65mm, 0) --
    ++(130mm, 0);

  % Featured article 2
  \node[align=left, text width=130mm, inner sep=0] at ([yshift=-134mm]current page.north) {
    {\sffamily\bfseries\color{accent2}\footnotesize\MakeUppercase{Perspective}}\\[1pt]
    {\sffamily\fontsize{16}{20}\selectfont\bfseries\color{accent}
     CRISPR 2.0: Precision Gene Editing Enters the Clinic}\\[3pt]
    {\sffamily\footnotesize\color{graytext}\itshape
     Sarah Nakamura\ \textbar\ page 9}
  };

  % Rule
  \draw[lightrule, line width=0.4pt]
    ([yshift=-147mm]current page.north) ++(-65mm, 0) --
    ++(130mm, 0);

  % Featured article 3
  \node[align=left, text width=130mm, inner sep=0] at ([yshift=-156mm]current page.north) {
    {\sffamily\bfseries\color{accent2}\footnotesize\MakeUppercase{Review}}\\[1pt]
    {\sffamily\fontsize{16}{20}\selectfont\bfseries\color{accent}
     Climate Models at the Exascale: Simulating Earth's Future}\\[3pt]
    {\sffamily\footnotesize\color{graytext}\itshape
     Marcus Okafor, Priya Sharma \& The CMIP7 Consortium\ \textbar\ page 15}
  };

  % Bottom publisher info
  \node[graytext, font=\sffamily\footnotesize, align=center] at ([yshift=14mm]current page.south) {
    Published by the Scientific Horizons Press \textbar\ Established 2012\\
    ISSN 2754-0182 \textbar\ www.scientifichorizons.quarterly
  };

\end{tikzpicture}

\restoregeometry
\pagecolor{pagebg}

% ====== PAGE 2 : MASTHEAD + EDITORIAL =======================================
\newpage
\thispagestyle{empty}

\begin{center}
{\sffamily\bfseries\fontsize{20}{24}\selectfont\color{accent}Scientific Horizons}

\vspace{6pt}
{\sffamily\color{graytext}\footnotesize
A Quarterly Review of Discovery \& Innovation \quad·\quad
Vol.\ 12, No.\ 3, Autumn 2024}

\vspace{4pt}
{\color{lightrule}\hrule height 0.5pt}
\end{center}

\vspace{10pt}

% Masthead using tabularx with l-column keys (NOT p{})
{\sffamily\footnotesize
\setlength{\tabcolsep}{0pt}
\begin{tabularx}{\linewidth}{
  @{}>{\color{accent}\bfseries}l@{\hspace{6mm}}X@{\hspace{10mm}}
  >{\color{accent}\bfseries}l@{\hspace{6mm}}X@{}
}
Editor-in-Chief     & Prof.\ Margaret Okonkwo, FRS
                     & Managing Editor    & Dr.\ Robert Lindqvist \\
Senior Editors      & Dr.\ Amara Singh,
                       Prof.\ Kenji Tanaka
                     & Reviews Editor     & Dr.\ Lucia Fernandez \\
Production          & Elena Petrova,
                       James Whelan
                     & Online Editor      & Dr.\ Samuel Greene \\
\end{tabularx}
}

\vspace{8pt}
{\color{lightrule}\hrule height 0.3pt}
\vspace{8pt}

{\sffamily\footnotesize\color{graytext}
\textbf{Editorial Board:} Prof.\ Alicia Montoya (Columbia),
Prof.\ David Okonkwo (Cambridge), Prof.\ Ravi Chandrasekhar (IISc),
Prof.\ Li Wei (Tsinghua), Prof.\ Elena Sokolova (Moscow State),
Prof.\ Carlos Ruiz (UNAM), Prof.\ Fatima Al-Rashid (KAUST),
Prof.\ Henrik Johansson (Karolinska), Prof.\ Yuki Nakamura (Tokyo),
Prof.\ Sarah Mitchell (Melbourne), Prof.\ Jean-Pierre Dubois (Sorbonne),
Prof.\ Kwame Asante (Cape Town).\par
}

\vspace{6pt}
{\sffamily\footnotesize\color{graytext}
\textbf{Contact:} Scientific Horizons Press, 42 Cavendish Square,
London W1G 0AN, United Kingdom. {\footnotesize\texttt{editorial.\allowbreak quarterly}}
\par}

\vspace{12pt}
{\color{lightrule}\hrule height 0.5pt}
\vspace{14pt}

% Editorial
{\sffamily\bfseries\Large\color{accent}From the Editor}

\vspace{8pt}

\dropcap{T}{he} practice of science has always been shaped by its tools. The
telescope, the microscope, the particle accelerator --- each opened a new window
onto nature and, in doing so, transformed not only what we know but how we
think. Today, we stand at the threshold of another such transformation, driven
not by a single instrument but by a confluence of computational power,
algorithmic ingenuity, and data abundance.

\medskip
In this issue of \textit{Scientific Horizons}, our feature article by Vasquez
and Chen examines one of the most exciting frontiers in this new landscape: the
application of machine learning to the discovery of quantum materials. The
promise is extraordinary --- algorithms that can predict stable crystal
structures, identify candidates for superconductivity, and even suggest
synthetic pathways, all before a single experiment is performed. But as our
authors remind us, the algorithmic tail must not wag the scientific dog. Close
collaboration between computational and experimental researchers remains
essential.

\medskip
Equally compelling is Nakamura's perspective on the clinical translation of
CRISPR-based therapies. A decade after the first demonstrations of programmable
gene editing in human cells, we are now seeing treatments move from bench to
bedside. The challenges --- delivery, specificity, and equitable access --- are
daunting, but the trajectory is unmistakable.

\medskip
Finally, our review by Okafor, Sharma, and the CMIP7 Consortium surveys the
state of climate modelling at the exascale. As computing power grows, so too
does our ability to resolve the fine-scale processes that govern cloud
formation, ocean mixing, and ice-sheet dynamics. The picture that emerges is
sobering but actionable --- a call to both scientific and policy communities.

\medskip
We hope this issue informs, challenges, and inspires. Science is a collective
endeavour, and every reader of these pages is part of the conversation.

\vspace{10pt}
{\sffamily\itshape\color{graytext}--- Margaret Okonkwo, Editor-in-Chief}

\vspace{8pt}
{\sffamily\footnotesize\color{graytext}London, September 2024}

% ====== PAGE 3 : TABLE OF CONTENTS ==========================================
\newpage
\thispagestyle{empty}

{\sffamily\bfseries\fontsize{22}{26}\selectfont\color{accent}Contents}

\vspace{6pt}
{\color{accent2}\hrule height 1.5pt}
\vspace{14pt}

% --- DEPARTMENTS ---
{\sffamily\bfseries\footnotesize\color{accent2}\MakeUppercase{Departments}}

\vspace{8pt}

{\sffamily
\setlength{\tabcolsep}{0pt}
\begin{tabularx}{\linewidth}{@{}>{\bfseries\color{accent}}l@{\hspace{6mm}}X@{\hspace{8mm}}r@{}}
2  & From the Editor \dotfill & 2 \\
4  & Research Highlights \dotfill & 4 \\
22 & Books \& Reviews \dotfill & 22 \\
24 & Calendar \& Announcements \dotfill & 24 \\
\end{tabularx}
}

\vspace{16pt}

% --- FEATURES ---
{\sffamily\bfseries\footnotesize\color{accent2}\MakeUppercase{Feature Articles}}

\vspace{8pt}

{\sffamily
\setlength{\tabcolsep}{0pt}
\begin{tabularx}{\linewidth}{@{}>{\bfseries\color{accent}}l@{\hspace{6mm}}X@{\hspace{8mm}}r@{}}
4  & The Computational Frontier: Machine Learning Meets Quantum
     Materials Discovery \dotfill & 4 \\
   & \hspace{2mm}\itshape\color{graytext}\footnotesize Elena Vasquez \& James Chen & \\
\end{tabularx}
}

\vspace{10pt}

{\sffamily
\setlength{\tabcolsep}{0pt}
\begin{tabularx}{\linewidth}{@{}>{\bfseries\color{accent}}l@{\hspace{6mm}}X@{\hspace{8mm}}r@{}}
9  & CRISPR 2.0: Precision Gene Editing Enters the Clinic \dotfill & 9 \\
   & \hspace{2mm}\itshape\color{graytext}\footnotesize Sarah Nakamura & \\
\end{tabularx}
}

\vspace{10pt}

{\sffamily
\setlength{\tabcolsep}{0pt}
\begin{tabularx}{\linewidth}{@{}>{\bfseries\color{accent}}l@{\hspace{6mm}}X@{\hspace{8mm}}r@{}}
15 & Climate Models at the Exascale: Simulating Earth's Future \dotfill & 15 \\
   & \hspace{2mm}\itshape\color{graytext}\footnotesize Marcus Okafor, Priya Sharma
      \& The CMIP7 Consortium & \\
\end{tabularx}
}

\vspace{20pt}
{\color{lightrule}\hrule height 0.5pt}
\vspace{10pt}

% Cover image credit
{\sffamily\footnotesize\color{graytext}
\textbf{Cover:} Crystallographic structure of the perovskite
CaTiO\textsubscript{3}, rendered from neutron diffraction data. Image courtesy
of the Materials Project (Lawrence Berkeley National Laboratory).\\
\textbf{Design:} Elena Petrova. Typeset in \TeX\ Gyre Pagella and Fira Sans.
}

% ====== PAGE 4-8 : FEATURE ARTICLE 1 ========================================
\newpage

% Article header (full width before columns)
{\sffamily\bfseries\color{accent2}\footnotesize\MakeUppercase{Feature Article}}

\vspace{4pt}
{\sffamily\bfseries\fontsize{24}{30}\selectfont\color{accent}
The Computational Frontier: Machine Learning Meets Quantum Materials Discovery}

\byline{Elena Vasquez\thanks{Department of Materials Science, Stanford University,
Stanford, CA 94305, USA.} \hspace{4pt}$\cdot$\hspace{4pt}
James Chen\thanks{Theory of Condensed Matter Group, Cavendish Laboratory,
University of Cambridge, Cambridge CB3 0HE, UK.}}

\begin{multicols}{2}

\dropcap{I}{n} the summer of 2023, a research group at Nexa DeepMind published
a landmark study in \textit{Nature} reporting the computational prediction of
over 380,000 stable crystalline materials --- a number that nearly doubled the
cumulative output of a century of experimental crystallography. The tool they
used, Graph Networks for Materials Exploration (GNoME), relied on deep
graph-neural networks trained on data from the Materials Project, a
US Department of Energy initiative that has systematically catalogued the
properties of known inorganic compounds since 2011.

\medskip
This was not an isolated breakthrough. Across condensed-matter physics and
materials science, machine-learning (ML) techniques are rapidly reshaping the
discovery pipeline. From predicting superconducting transition temperatures to
identifying optimal dopants for battery cathodes, data-driven approaches are
increasingly complementing --- and in some cases supplanting --- traditional
methods rooted in density functional theory (DFT) and empirical intuition.

\medskip
Yet the convergence of ML and materials science raises fundamental questions
that extend beyond raw predictive power. How do we validate
computationally-predicted materials when the volume of predictions outstrips the
capacity of experimental verification by orders of magnitude? What does it mean
for a crystal structure to be ``stable'' when the definition depends on the
choice of reference chemical potentials? And critically, can ML models ever
generate genuine physical insight, or are they merely sophisticated
interpolators operating within the boundaries of their training data?

\medskip
The present article surveys the state of this rapidly evolving field and
identifies three grand challenges that will determine whether the marriage of
ML and materials science fulfills its considerable promise.

\subsubsection{From DFT to Deep Learning}

The computational study of materials has, for four decades, been dominated by
DFT and the family of methods built upon it. DFT-based high-throughput
screening --- in which thousands of candidate structures are evaluated against
thermodynamic stability criteria --- has become a standard workflow, exemplified
by the Materials Project, AFLOW, and the Open Quantum Materials Database
(OQMD). These databases now contain over two million entries and have enabled
the discovery of numerous functional materials, including new
thermoelectrics, transparent conductors, and Li-ion conductors.

\medskip
However, DFT-based screening faces inherent limitations. First-principles
calculations are computationally expensive, scaling as $O(N^3)$ with system
size, making exhaustive exploration of compositional and configurational space
infeasible. Moreover, standard exchange-correlation functionals (e.g., PBE)
introduce systematic errors in band gaps, formation energies, and
reaction barriers --- quantities central to materials design. Hybrid
functionals and many-body methods (GW, RPA) improve accuracy but at prohibitive
computational cost.

\medskip
ML models offer a compelling alternative by learning the mapping from structure
to property directly from data. Rather than solving the Kohn--Sham equations
for each candidate, a trained model can evaluate stability, band structure, or
mechanical properties in milliseconds. Early work focused on kernel-based
methods (Gaussian processes, support vector regression) trained on
hand-crafted descriptors. The field has since shifted decisively toward
deep-learning architectures, which can learn representations directly from
atomic coordinates and species.

\subsubsection{Representation and Architecture}

Three classes of neural-network architecture now dominate the materials
informatics landscape. Crystal graph convolutional networks (CGCNs) represent
a crystal as a graph in which nodes correspond to atoms and edges encode
interatomic distances. Message-passing layers aggregate information from
neighbouring atoms, recursively building a representation of the local chemical
environment. The MatErials Graph Network (MEGNet) and Crystal Graph
Convolutional Neural Network (CGCNN) frameworks, both introduced in
2018--2019, remain among the most widely used.

\medskip
Equivariant neural networks, which explicitly respect the symmetry group of
three-dimensional space (E(3)), represent a more recent advance. By
constraining their internal representations to transform correctly under
rotations and translations, these models achieve state-of-the-art accuracy with
fewer training examples. The NequIP and MACE architectures, built on
higher-order tensor products of spherical harmonics, have demonstrated
remarkable performance on the prediction of interatomic forces and phonon
spectra.

\medskip
Finally, diffusion models and flow-matching approaches, originally developed
for image and video generation, are now being adapted to generate crystal
structures \textit{de novo}. These generative models learn to reverse a
noise-adding process, producing candidate structures that satisfy specified
property constraints. GNoME combined a graph-network stability predictor with an
active-learning loop, iteratively proposing candidates, evaluating them with
DFT, and retraining on the expanded dataset.

\medskip
The practical impact of these methods is already evident. A 2024 study by
Merchant \textit{et al.} used GNoME to identify 52 previously unknown layered
materials with potential as solid electrolytes. Independent DFT validation
confirmed the stability of 48 of these candidates, and experimental synthesis
has been initiated at three laboratories. Similarly, Chen \textit{et al.}
reported a graph-network-driven search for high-entropy alloy catalysts that
yielded two compositions outperforming commercial Pt/C in oxygen reduction
reaction activity.

\pullquote{The algorithmic tail must not wag the scientific dog. Close
collaboration between computational and experimental researchers remains
essential for realizing the promise of ML-driven discovery.}

\subsubsection{Three Grand Challenges}

Despite these successes, the field confronts obstacles that will require
concerted effort across the computational and experimental communities to
overcome.

\textbf{Challenge 1: Validation at scale.} When a single ML study can predict
hundreds of thousands of stable compounds, the bottleneck shifts decisively to
experiment. Automated synthesis platforms --- including robotic solid-state
synthesis and high-throughput thin-film deposition --- are beginning to address
this gap, but throughput remains orders of magnitude below what computational
prediction can produce. The community must develop tiered validation
protocols that prioritise candidates with the highest probability of successful
synthesis and the greatest potential impact.

\textbf{Challenge 2: Uncertainty quantification.} Most ML models for materials
prediction provide point estimates without meaningful confidence intervals.
When a model predicts a formation energy of $-1.2$\,eV/atom, how certain are we
that the compound is stable? Bayesian neural networks and ensemble methods
offer partial solutions, but rigorous uncertainty quantification (UQ) tailored
to materials-specific error modes --- configurational entropy, metastable
polymorphs, anharmonic effects --- remains underdeveloped.

\textbf{Challenge 3: From prediction to understanding.} The most profound
critique of ML in the physical sciences is that it trades understanding for
prediction. A model that accurately forecasts superconducting $T_c$ without
revealing the underlying pairing mechanism may be useful for engineering but
contributes little to fundamental physics. Emerging interpretability techniques
--- including layer-wise relevance propagation, concept activation vectors, and
symbolic regression on learned representations --- aim to extract
human-comprehensible insights from trained networks, but the gap between these
tools and the depth of understanding afforded by traditional theory remains
wide.

% Infographic callout box
\begin{infobox}
{\sffamily\bfseries\color{accent}\small\MakeUppercase{By the Numbers}}

\vspace{4pt}
{\sffamily\footnotesize
\setlength{\tabcolsep}{0pt}
\begin{tabularx}{\linewidth}{@{}>{\bfseries\color{accent2}}lX@{}}
380K  & Stable materials predicted by GNoME (2023) \\
200K  & Known inorganic compounds in ICSD (2024) \\
2.2M  & Total entries across Materials Project, AFLOW, OQMD \\
10\textsuperscript{4} & Speedup of ML inference vs.\ hybrid DFT \\
48/52 & GNoME-predicted Li conductors validated by DFT \\
\end{tabularx}
}
\end{infobox}

\subsubsection{The Road Ahead}

The convergence of ML and materials science is not merely a technological trend
but a fundamental shift in how we explore chemical and structural space. The
field is evolving from a paradigm of \textit{predict-then-validate} toward one
of \textit{co-design}, in which computational prediction, automated synthesis,
and characterisation are integrated into closed-loop workflows. Early
demonstrations of autonomous laboratories --- such as the A-Lab at Lawrence
Berkeley National Laboratory, which combines ML-driven experiment planning with
robotic synthesis and X-ray diffraction characterisation --- point toward a
future in which the discovery cycle is compressed from years to days.

\medskip
Realising this vision will require not only algorithmic advances but also
cultural change. Materials scientists and condensed-matter physicists must
become fluent in the language of machine learning, while computer scientists
must engage seriously with the physics that makes materials prediction a
uniquely challenging domain. The journals, funding agencies, and academic
departments that structure scientific careers must adapt to reward
interdisciplinary work that bridges these historically separate communities.

\medskip
The stakes are high. From better batteries and more efficient solar cells to
room-temperature superconductors and quantum information platforms, the
materials we need to address the twenty-first century's most pressing challenges
are waiting to be discovered. Machine learning, deployed thoughtfully and in
partnership with experiment and theory, may be the tool that helps us find them.

\vspace{8pt}
{\sffamily\footnotesize\color{graytext}
\textbf{Acknowledgements.} E.V.\ acknowledges support from the US Department of
Energy, Office of Basic Energy Sciences (DE-SC0022417). J.C.\ acknowledges
support from the European Research Council (Grant No.\ 101054458) and the
Leverhulme Trust. The authors thank A.\ Merchant, R.\ Gomez-Bombarelli, and
K.\ Persson for illuminating discussions.
}

\vspace{4pt}
{\sffamily\footnotesize\color{graytext}
\textbf{Competing interests.} The authors declare no competing interests.
}

\vspace{6pt}
{\sffamily\footnotesize\color{graytext}
\textbf{Data availability.} All datasets referenced in this article are publicly
available through the Materials Project (materialsproject.org), ICSD, and
associated repositories.
}

\end{multicols}

% ====== PAGE 9 : FEATURE ARTICLE 2 ==========================================
\newpage

{\sffamily\bfseries\color{accent2}\footnotesize\MakeUppercase{Perspective}}

\vspace{4pt}
{\sffamily\bfseries\fontsize{24}{30}\selectfont\color{accent}
CRISPR 2.0: Precision Gene Editing Enters the Clinic}

\byline{Sarah Nakamura\thanks{Division of Genetics, Brigham and Women's Hospital,
Harvard Medical School, Boston, MA 02115, USA.}}

\begin{multicols}{2}

\dropcap{W}{hen} Victoria Gray became the first person to receive an
investigational CRISPR-based therapy for sickle-cell disease in 2019, the
procedure was rightly hailed as a landmark. Four years later, in December 2023,
the US Food and Drug Administration approved Casgevy (exagamglogene autotemcel),
the first CRISPR/Cas9-based therapeutic, for both sickle-cell disease and
transfusion-dependent beta-thalassemia. The approval, which followed similar
authorisation by the UK Medicines and Healthcare products Regulatory Agency,
marked the culmination of a journey that began with the discovery of CRISPR
adaptive immunity in bacteria and the demonstration, by Doudna and Charpentier
in 2012, that the Cas9 nuclease could be programmed for sequence-specific DNA
cleavage.

\medskip
Yet the approval of Casgevy, momentous as it is, represents only the first
chapter in the clinical translation of genome editing. A second generation of
technologies --- base editing, prime editing, and epigenetic editing --- is
already entering human trials, promising greater precision, broader
applicability, and fewer off-target effects than first-generation CRISPR-Cas9.

\subsubsection{Beyond Double-Strand Breaks}

First-generation CRISPR-Cas9 therapy relies on the introduction of a
double-strand break (DSB) at a targeted genomic locus, followed by repair
through either non-homologous end joining (NHEJ) or homology-directed repair
(HDR). While effective --- Casgevy works by disrupting the BCL11A erythroid
enhancer to reactivate fetal haemoglobin production --- DSB-based editing
carries inherent risks, including large deletions, complex rearrangements, and
activation of the p53 DNA-damage response.

\medskip
Base editing, pioneered by David Liu's laboratory at the Broad Institute,
sidesteps DSBs entirely. Cytosine base editors (CBEs) and adenine base editors
(ABEs) fuse a catalytically-impaired Cas9 nickase to a deaminase enzyme,
enabling the direct conversion of C·G to T·A or A·T to G·C base pairs without
cleaving the DNA backbone. The first clinical trial of a base editor,
BEAM-101 for sickle-cell disease, dosed its first patient in 2023. Early
results, presented at the 2024 American Society of Hematology meeting, showed
robust HbF induction and no evidence of off-target mutations by whole-genome
sequencing.

\medskip
Prime editing, also developed in Liu's laboratory, extends the editing
repertoire further. By fusing a Cas9 nickase to an engineered reverse
transcriptase and providing a prime editing guide RNA (pegRNA) that both
specifies the target and encodes the desired edit, prime editors can perform all
twelve possible base-to-base conversions, as well as small insertions (up to
\~{}44\,bp) and deletions (up to \~{}80\,bp), without requiring DSBs or donor
templates. The first IND application for a prime-editing therapy --- targeting
chronic granulomatous disease --- is anticipated in 2025.

\pullquote{The approval of Casgevy represents only the first chapter in the
clinical translation of genome editing. A second generation of technologies
is already entering human trials.}

\subsubsection{Delivery Remains the Bottleneck}

For all the sophistication of editing modalities, delivery remains the central
challenge in translating genome-editing therapies to the clinic. The \textit{ex
vivo} approach used by Casgevy --- harvesting haematopoietic stem cells,
editing them in culture, and reinfusing after myeloablative conditioning --- is
effective but expensive (Casgevy is priced at \$2.2 million per patient in the
US) and logistically demanding, requiring specialised manufacturing facilities
and extended hospitalisation.

\medskip
\textit{In vivo} delivery, in which the editing machinery is administered
directly to the patient, would dramatically expand the addressable patient
population. Lipid nanoparticles (LNPs), the delivery vehicles that enabled the
success of mRNA COVID-19 vaccines, have shown promise for hepatic delivery of
CRISPR components. Intellia Therapeutics' NTLA-2001, an LNP-formulated
Cas9 mRNA and guide RNA targeting the \textit{TTR} gene for transthyretin
amyloidosis, demonstrated durable (\textgreater 90\%) TTR knockdown in Phase
I/II trials, representing the first clinical proof-of-concept for systemic
\textit{in vivo} CRISPR editing.

\medskip
For tissues beyond the liver, virus-like particles (VLPs) and engineered
adeno-associated viruses (AAVs) are under active investigation, though
immunogenicity and payload-size constraints remain significant hurdles.

\subsubsection{The Ethical Horizon}

As genome-editing technologies become more precise and more accessible, the
ethical questions they raise grow sharper. Somatic editing for severe genetic
diseases commands broad consensus; germline editing, which would introduce
heritable changes, does not. The 2018 announcement by He Jiankui of the birth of
genome-edited twins, universally condemned by the scientific community, served
as a stark reminder that technological capability can outpace ethical consensus.

\medskip
The World Health Organization's 2021 governance framework for human genome
editing, the International Commission on the Clinical Use of Human Germline
Genome Editing's 2020 report, and the 2023 Third International Summit on Human
Genome Editing in London have each contributed to a gradually coalescing global
framework. Whether that framework can accommodate the accelerating pace of
technological development --- and whether it can ensure equitable access to
therapies that could otherwise deepen health disparities --- remains to be seen.

\vspace{8pt}
{\sffamily\footnotesize\color{graytext}
\textbf{Acknowledgements.} S.N.\ is supported by NIH grant R01-HG012034 and the
Harvard Stem Cell Institute. She serves on the scientific advisory board of
Verve Therapeutics (uncompensated).
}

\vspace{4pt}
{\sffamily\footnotesize\color{graytext}
\textbf{Competing interests.} S.N.\ holds equity in Beam Therapeutics and is a
named inventor on patents related to base-editing technologies licensed to Beam
Therapeutics through the Broad Institute.
}

\end{multicols}

% ====== PAGE 11 : RESEARCH HIGHLIGHTS =======================================
\newpage

\section{Research Highlights}

\subsection*{Astrophysics}

{\sffamily\bfseries\color{accent}JWST Reveals Carbonaceous Dust in the Early
Universe.} Observations from the James Webb Space Telescope have detected the
spectral signature of polycyclic aromatic hydrocarbons (PAHs) in a galaxy at
$z = 8.3$, corresponding to an epoch just 600 million years after the Big Bang.
The finding, reported by Witstok \textit{et al.} in \textit{Nature Astronomy}
(2024), challenges existing models of dust production, which require
intermediate-mass AGB stars that would not have had time to evolve by $z > 8$.
The authors propose that carbonaceous dust may form in the ejecta of
core-collapse supernovae on shorter timescales than previously recognised, a
hypothesis consistent with recent nucleosynthetic models. \textit{--- Nature
Astronomy} {\bf 8}, 1124--1131 (2024)

\subsection*{Neuroscience}

{\sffamily\bfseries\color{accent}A Molecular Atlas of the Human Blood--Brain
Barrier.} A multi-institutional consortium led by the Allen Institute for Brain
Science has published a comprehensive single-cell transcriptomic atlas of the
human blood--brain barrier (BBB), profiling over 140,000 cells from 17 brain
regions across 8 donors. The atlas, described in \textit{Cell} (2024), reveals
previously unrecognised heterogeneity in endothelial transport machinery,
identifies region-specific vulnerabilities relevant to neurodegenerative
disease, and provides a reference dataset for efforts to engineer the BBB for
therapeutic delivery. \textit{--- Cell} {\bf 187}, 3487--3504 (2024)

\subsection*{Climate Science}

{\sffamily\bfseries\color{accent}Accelerating Ice Loss from the Amundsen Sea
Sector.} A new analysis combining satellite altimetry, ice-penetrating radar,
and oceanographic mooring data indicates that ice loss from West Antarctica's
Amundsen Sea sector has accelerated by 35\% over the past decade. The study,
published in \textit{Science Advances} by Davison \textit{et al.} (2024),
attributes the acceleration to increased intrusions of warm Circumpolar Deep
Water onto the continental shelf, driven by shifts in the Amundsen Sea Low
associated with anthropogenic forcing. The findings raise the estimated
contribution of the Amundsen sector to global sea-level rise to
0.6\,mm\,yr$^{-1}$. \textit{--- Sci.\ Adv.} {\bf 10}, eadk7890 (2024)

\subsection*{Synthetic Biology}

{\sffamily\bfseries\color{accent}De Novo Design of Allosteric Protein Switches.}
Researchers at the University of Washington's Institute for Protein Design have
used deep-learning-based protein design (RFdiffusion and ProteinMPNN) to create
synthetic allosteric switches --- proteins that change conformation and activity
in response to binding a small-molecule ligand --- entirely from first
principles. The designed switches, reported in \textit{Science} by Anishchenko
\textit{et al.} (2024), achieve switching ratios exceeding 50-fold and
demonstrate the feasibility of computational design of conditional protein
function. \textit{--- Science} {\bf 385}, 178--185 (2024)

\vspace{10pt}
{\color{lightrule}\hrule height 0.3pt}

% ====== CONTRIBUTOR BIOS ====================================================
\vspace{14pt}
\section{Contributors}

\begin{multicols}{2}

\contributor{Elena Vasquez}
{Assistant Professor, Department of Materials Science, Stanford University}
{Elena Vasquez leads the Computational Materials Design group at Stanford. Her
research combines first-principles electronic-structure calculations with
machine-learning methods to accelerate the discovery of functional materials for
energy applications. She received her PhD from MIT in 2018 and was a Miller
Research Fellow at UC Berkeley before joining the Stanford faculty in 2021.}

\contributor{James Chen}
{Professor of Theoretical Physics, Cavendish Laboratory, University of
Cambridge}
{James Chen holds the Cavendish Chair in Condensed Matter Theory. His group
develops and applies machine-learning interatomic potentials and equivariant
neural networks to problems in materials physics, with a particular focus on
ferroelectrics, multiferroics, and quantum materials. He is a Fellow of the
Royal Society and a recipient of the Dirac Medal.}

\contributor{Sarah Nakamura}
{Assistant Professor, Division of Genetics, Brigham and Women's Hospital,
Harvard Medical School}
{Sarah Nakamura is a physician-scientist whose research spans the development of
novel genome-editing tools and their translation to clinical application. Her
laboratory pioneered several base-editing strategies for haematological
disorders. She trained in David Liu's laboratory at the Broad Institute and
completed her clinical fellowship in medical genetics at Boston Children's
Hospital.}

\contributor{Marcus Okafor}
{Senior Research Scientist, Met Office Hadley Centre, Exeter}
{Marcus Okafor leads the High-Resolution Climate Modelling group at the Met
Office. His research focuses on the representation of cloud and convective
processes in global climate models, with an emphasis on leveraging exascale
computing to resolve kilometre-scale atmospheric dynamics. He serves on the
CMIP7 Panel and the WCRP Grand Challenge on Clouds and Circulation.}

\contributor{Priya Sharma}
{Associate Professor, Department of Earth System Science, University of
California, Irvine}
{Priya Sharma's group develops and applies coupled Earth-system models,
integrating atmospheric, oceanic, land-surface, and ice-sheet components. Her
current research focuses on ice-sheet--ocean interactions in Antarctica and
their implications for sea-level projections under high-emission scenarios. She
is a lead author on the IPCC Seventh Assessment Report.}

\end{multicols}

\vspace{10pt}
{\color{lightrule}\hrule height 0.5pt}
\vspace{10pt}

\begin{center}
{\sffamily\footnotesize\color{graytext}
\textit{Scientific Horizons} is published quarterly by Scientific Horizons Press,\\
42 Cavendish Square, London W1G 0AN, United Kingdom.\\
ISSN\ 2754-0182 (Print) \quad·\quad ISSN\ 2754-0190 (Online)\\[4pt]
\textcopyright\ 2024 Scientific Horizons Press. All rights reserved.\\
No part of this publication may be reproduced without prior written permission.\\[4pt]
Typeset in \LaTeX\ using \TeX\ Gyre Pagella and Fira Sans.\\
Printed on FSC-certified paper by Clays Ltd, Suffolk.
}
\end{center}

\end{document}
